%! Equivalent Representations of Polynomial and Rational Expressions — Unit 3, Lesson 3.5 — Student Packet Cover Sheet
\documentclass[10pt]{article}
\usepackage{precalculus-article}
\usepackage{precalculus-boxes}
\usepackage{ltablex}
\keepXColumns

\begin{document}

% Full-bleed navy banner
\begin{tikzpicture}[remember picture, overlay]
  \fill[navy] (current page.north west)
    rectangle ([yshift=-0.9in]current page.north east);
\end{tikzpicture}
\vspace{-0.8in}

\begin{center}
  {\color{white}\textbf{\LARGE Precalculus}}\\[4pt]
  {\color{skymid}\large\textbf{Unit 3: Rational Functions}}\\[6pt]
  {\color{white}\textbf{\large Lesson 3.5 \quad Equivalent Representations of Polynomial and Rational Expressions}}
\end{center}

\vspace{0.22in}
\namedateperiod
\vspace{0.18in}

\begin{learningtargetbox}
\small
\begin{enumerate}[leftmargin=1.5em, itemsep=5pt]
  \item I can rewrite a polynomial or rational expression in factored or standard form and
        explain what each form reveals about the function (zeros, end behavior, asymptotes, holes).
  \item I can perform polynomial long division to express $f(x) = g(x)q(x) + r(x)$, identify a
        slant asymptote from the quotient, and use Pascal's Triangle to expand $(x+c)^n$.
\end{enumerate}
\end{learningtargetbox}

\vspace{0.14in}

\begin{tocbox}
\small
\renewcommand{\arraystretch}{1.5}
\begin{tabularx}{\linewidth}{@{} c l X r @{}}
\rowcolor{sky}
\textbf{\#} & \textbf{Component} & \textbf{Description} & \textbf{Score} \\
\midrule
1 & Warm-Up        & Spiral review: factoring, rational zeros, asymptotes, and holes & \blank{1.2cm} \\
2 & Guided Notes   & Factored vs.\ standard form; polynomial long division; binomial theorem & \blank{1.2cm} \\
3 & Group Activity & Form-switching with rational functions (tiered R / A / E)         & \blank{1.2cm} \\
4 & Exit Ticket    & Choose a form; execute one long-division step; Pascal expansion    & \blank{1.2cm} \\
5 & Homework       & Mixed: forms, long division, slant asymptotes, binomial theorem    & \blank{1.2cm} \\
\midrule
  & \multicolumn{2}{r}{\textbf{Total}} & \blank{1.2cm} \\
\end{tabularx}
\end{tocbox}

\end{document}
